\begin{ledgroupsized}[r]{120mm}
\footnotesize
\pstart
\noindent\textbf{Überlieferung:}
\pend
\end{ledgroupsized}
\begin{ledgroupsized}[r]{114mm}
\footnotesize
\pstart
\parindent -6mm
\makebox[6mm][l]{\textit{L}}%
Konzept:
LH~XXXV~11, 14~Bl.~8.
Ein Blatt~4\textsuperscript{o};
\href{\Scan{LH_35_11_14_008wz}}%
{Wasserzeichen} in der Blattmitte;
alle Ränder bis auf den oberen beschnitten.
Eineinhalb Seiten.
\pend
\end{ledgroupsized}
\vspace{5mm}
\begin{ledgroup}
\footnotesize
\pstart
\noindent
\textbf{Datierungsgründe:}
Das Wasserzeichen %
ist mehrfach für den angegebenen Zeitraum belegt:
einerseits in einigen in \indexO{Hannover}Hannover geschriebenen Briefen 
\indexN{\textls{Guidi} Giuseppe gest. 1720\,(?)}Giuseppe Guidis
(Briefe vom 3. und 7. November 1703: \textit{LSB}~I, 22 \cite{02156}N.\;381 und \cite{02157}N.\;387, sowie \cite{02158}Brief von April 1704: \textit{LSB}~I, 23 N.\;226);
andererseits im Träger des eigenhändig auf 1708 datierten \glqq Entwurff einer Historischen Medaille auf die Introductionem in das Collegium Electorale\grqq\  %
(LH XXIII 2, 18 Bl.~4; erscheint vsl. in einem Band der Vierten Reihe).
Aus diesem Grund wird von einer Entstehung von N.\;\ref{58942} in der angegebenen Zeitspanne ausgegangen.
\pend 
\end{ledgroup}
\selectlanguage{latin}
\vspace{8mm}
\pstart
\noindent
\href{\Scan{LH_35_11_14_008r}}%
{\lbrack8~r\textsuperscript{o}\rbrack}\ 
Centrum Gravitatis\indexS{centrum gravitatis}
\edtext{ubique}{\lemma{}\Bfootnote{ubique \textit{erg.~L}}} 
existere\indexS{demonstratio quod centrum gravitatis ubique existit} ex hoc principio demonstro, quod datis duobus circulis 
\edtext{vel aliis figuris\indexS{figura centrum gravitatis habens} centrum gravitatis\indexS{centrum gravitatis} haud dubie habentibus}{\lemma{}\Bfootnote{vel aliis figuris centrum gravitatis haud dubie habentibus \textit{erg.~L}}} 
datur punctum, per quod ducta quaevis recta facit ut circuli duo sint in aequilibrio\indexS{aequilibrium} seu ut momentum\indexS{momentum}\indexS{bisectio momenti in centro gravitatis} bisecetur\lbrack:\rbrack\ 
\edtext{hinc resoluta}{\lemma{hinc}\Bfootnote{%
\textit{(1)}~resolutio %
\textit{(2)}~resoluta~\textit{L}}}
Figura\indexS{resolutio figurae in circulos} in
\edtext{quotcunque circulos}{\lemma{quotcunque}\Bfootnote{%
\textit{(1)}~puncta %
\textit{(2)}~circulos~\textit{L}}}
inaequales\lbrack,\rbrack\ licet pro duobus semper 
\edtext{substitu\lbrack i\rbrack}{\lemma{}\Bfootnote{substituit \textit{L ändert Hrsg.}}} 
unus salvo aequilibrio,\indexS{aequilibrium}\indexS{substitutio salvo aequilibrio}
ipsis aequalis circa centrum gravitatis\indexS{centrum gravitatis} descriptus. 
\pend
\pstart
Si reperiri 
\edtext{posset punctum}{\lemma{posset}\Bfootnote{%
\textit{(1)}~recta
\textit{(a)}~bi
\textit{(b)}~bisecans %
\textit{(2)}~punctum~\textit{L}}}
per quod recta quaevis ducta utrinque 
\edtext{aequalem quantitatem ex duobus circulis haberet; dari posset}{\lemma{aequalem}\Bfootnote{%
\textit{(1)}~relinqueret %
\textit{(2)}~quantitatem \bars{ex duobus circulis \textit{erg.}} haberet; dari
\textit{(a)}~possent %
\textit{(b)}~posset~\textit{L}}}
centrum magnitudinis.\indexS{centrum magnitudinis}
\pend
\pstart
Duas quantitates centrum gravitatis\indexS{centrum gravitatis} 
\edtext{habentes (verb. gratia}{%
\lemma{habentes}%
\Bfootnote{%
\textit{(1)}~(
\textit{(2)}~(verb. gratia~\textit{L}}}
duos circulos) habere centrum gravitatis commune\indexS{centrum gravitatis commune} sic demonstro.\indexS{demonstratio quod centrum gravitatis commune existit}
\pend
\begin{edfigure}[1em][1em]
\includegraphics[scale=1]{%
gesamttex/edit_VIII-4/images/LH_35_11_14_008_d1_008r.pdf}
\edcaption{\edtext{\lbrack\textit{Fig.~1}\rbrack}{%
\lemma{\lbrack\textit{Fig.~1}\rbrack}%
\Cfootnote[nonum]{Zwei gestrichene Entwürfe zum Diagramm werden nicht wiedergegeben.}}}
\end{edfigure}
\pstart 
Primum datur centrum gravitatis\indexS{centrum gravitatis} circuli\indexS{centrum gravitatis circuli} 
\edtext{\textit{A}}{\lemma{}\Bfootnote{\textit{A} \textit{erg.~L}}}
id est punctum per quod quaevis recta in aequalia circulum momenta\indexS{momentum}\indexS{sectio in aequalia momenta} secat\lbrack;\rbrack\ 
demonstratum etiam est, posito centro gravitatis,\indexS{centrum gravitatis} verum esse ut momentum\indexS{momentum} corporis 
\edtext{sit ductum corporis in centri gravitatis distantiam.\indexS{ductus corporis in distantiam centri gravitatis} Sint circuli centris}{\lemma{sit}\Bfootnote{%
\textit{(1)}~distantia %
\textit{(2)}~ductum corporis in %
\textit{(a)}~dis %
\textit{(b)}~centri gravitatis distantiam. %
\textit{(aa)}~Sint %
\textit{(aaa)}~centra
\textit{(bbb)}~circuli %
\textit{(bb)}~Sint circuli centris~\textit{L}}}
\textit{A} et \textit{B}. Jungatur recta \textit{AB} et 
\edtext{secetur reciproce}{%
\lemma{secetur}%
\Bfootnote{%
\textit{(1)}~\bars{in \textit{streicht Hrsg.}}
\textit{(2)}~reciproce~\textit{L}}}
ad circulos in \textit{C}; ita ut sit \textit{AC}, ad \textit{BC} ut circulus \textit{B} ad circulum \textit{A}. Dico \textit{C} esse centrum duorum circulorum,\indexS{centrum gravitatis duorum} 
\edtext{seu punctum per quod quaevis \lbrack recta\rbrack\ ducta}{\lemma{seu}\Bfootnote{%
\textit{(1)}~rectam \bars{per \textit{streicht Hrsg.}} quam quaevis %
\textit{(2)}~punctum per quod \bars{quod \textit{streicht Hrsg.}} quaevis 
\bars{recta \textit{erg. Hrsg.}} ducta~\textit{L}}}
ut \textit{DE}, circulorum momentum\indexS{bisectio momenti in centro gravitatis} 
\edtext{bisecat. \textit{A}}{\lemma{bisecat.}\Bfootnote{%
\textit{(1)}~Ex %
\textit{(2)}~\textit{A}~\textit{L}}}
et \textit{B} agantur in \textit{DE} normales \textit{AF}, \textit{BG}, habemus triangula similia \textit{AFC}, \textit{BGC} ergo \textit{AF}, ad \textit{BG} ut \textit{AC} ad \textit{BC}, id est \textit{AF} ad \textit{BG} ut circulus \textit{B} ad 
\edtext{circulum \textit{A}, seu}{\lemma{circulum \textit{A},}\Bfootnote{%
\textit{(1)}~id est %
\textit{(2)}~seu~\textit{L}}}
reciproce ut distantiae centrorum. Ergo momenta\indexS{momenta aequalia} sunt aequalia. 
Et idem foret, si recta \textit{DE}, secaret ipsas magnitudines, neque enim amplius magnitudines sed centra spectamus. 
Hinc sequitur rursus quotcunque magnitudinum centrum habentium\indexS{magnitudo centrum gravitatis habens} dari centrum commune\indexS{centrum gravitatis commune quotcunque magnitudinum} 
\href{\Scan{LH_35_11_14_008v}}%
{\lbrack8~v\textsuperscript{o}\rbrack}  
\edtext{quia ut ex}{\lemma{quia}\Bfootnote{%
\textit{(1)}~duo %
\textit{(2)}~ex quotcunque unum %
\textit{(3)}~ut ex~\textit{L}}}
duobus fecimus, unum, ita ex composito ex duobus, et uno faciemus rursus unum\lbrack,\rbrack\ et ita porro.
Eodem argumento inveniemus centrum quasi-gravitatis,\indexS{centrum quasi-gravitatis} si 
\edtext{non distantiae,}{\lemma{non}\Bfootnote{%
\textit{(1)}~lineae %
\textit{(2)}~distantiae, \textit{L}}}
sed quadrata vel cubi distantiarum\indexS{quadratum vel cubus distantiae a centro gravitatis} sint 
\edtext{reciproce, ut quantitates. V.\,g.}{%
\lemma{reciproce, ut}%
\Bfootnote{%
\textit{(1)}~corpora
\textit{(2)}~quantitates.
\textit{(a)}~(
\textit{(b)}~V.\,g.~\textit{L}}}
ut sit \textit{A} ad \textit{B}, ut \textit{AC} qu. ad \textit{BC} qu. Sit enim $\surd A$ ad $\surd B$, ut \textit{AC} ad \textit{BC}. 
\edlabel{35_11_14_008_1a}%
\edtext{}{\xxref{35_11_14_008_1a}{35_11_14_008_1b}\lemma{}\Afootnote{\textit{Am Rand}: \Denarius}}%
Ergo fingendum tantum pro magnitudinibus 
\edtext{substitui quae}{\lemma{substitui}\Bfootnote{%
\textit{(1)}~earum %
\textit{(2)}~similes in ea \bars{ratione \textit{streicht Hrsg.}} %
\textit{(3)}~quae~\textit{L}}}
sint ut earum
\edtext{latera, fingendo quantitates}{\lemma{latera,}\Bfootnote{%
\textit{(1)}~servata %
\textit{(2)}~fingendo %
\textit{(a)}~corpora habere crassitiem %
\textit{(b)}~quantitates~\textit{L}}}
esse planas et habere crassitiem, eamque supra et infra augeri vel minui.\edlabel{35_11_14_008_1b}
\pend
\begin{edfigure}[1em][1em]
\includegraphics[scale=1]{%
gesamttex/edit_VIII-4/images/LH_35_11_14_008_d2_008v.pdf}
\edcaption{\lbrack\textit{Fig.~2}\rbrack}
\end{edfigure}
\pstart 
$\displaystyle\frac{A\cdot NO}{B\cdot RS}= \surd\displaystyle\frac{A\cdot LM}{B\cdot PQ}$.
\pend
